% ═══════════════════════════════════════════════════════════════════════════
% FRAGEKARTEN für DS 10 - Beobachtungsstunde
% Zum Ausdrucken, Ausschneiden und Laminieren
% ═══════════════════════════════════════════════════════════════════════════

\documentclass[11pt, a4paper]{article}

\usepackage[utf8]{inputenc}
\usepackage[T1]{fontenc}
\usepackage[ngerman]{babel}
\usepackage{helvet}
\renewcommand{\familydefault}{\sfdefault}

\usepackage[margin=1cm]{geometry}
\usepackage[table]{xcolor}
\usepackage{tikz}
\usepackage{array}
\usepackage{tabularx}
\usepackage{amssymb}

% Farben
\definecolor{spanisch}{HTML}{c41e3a}
\definecolor{grau}{HTML}{666666}
\definecolor{hellgrau}{HTML}{f5f5f5}
\definecolor{gruppeA}{HTML}{2a7a4b}  % Grün für A
\definecolor{gruppeB}{HTML}{1565c0}  % Blau für B
\definecolor{warning}{HTML}{b35c00}   % Orange für Warnungen

% Kein Seitenumbruch zwischen Karten
\usepackage{needspace}

% Spanisch-Befehl
\newcommand{\es}[1]{\textcolor{spanisch}{\textit{#1}}}

\pagestyle{empty}

\begin{document}

% ═══════════════════════════════════════════════════════════════════════════
% FRAGEKARTEN A (EINFACH) - 6 Stück pro Seite
% ═══════════════════════════════════════════════════════════════════════════

\begin{center}
{\LARGE\bfseries Fragekarten A -- Nivel A2}\\[0.3cm]
{\small 6 Exemplare zum Ausschneiden}
\end{center}

\vspace{0.5cm}

% Karte A - Template (wird 6x wiederholt)
\newcommand{\karteA}{%
\begin{tikzpicture}
\draw[thick, gruppeA, rounded corners=3pt] (0,0) rectangle (8.5,6.2);
\fill[gruppeA] (0,5.5) rectangle (8.5,6.2);
\node[white, font=\bfseries\large] at (4.25,5.85) {TARJETA A -- Nivel A2};

% Inhalt
\node[anchor=north west, text width=7.8cm] at (0.3,5.3) {
\textbf{1. Nombre:}\\[0.1cm]
\es{¿Cómo te llamas?}\\
{\scriptsize\color{grau} (Wie heißt du?)}

\vspace{0.3cm}

\textbf{2. Carácter:}\\[0.1cm]
\es{¿Cómo eres?}\\
{\scriptsize\color{grau} (Wie bist du? $\rightarrow$ simpático/a, inteligente...)}

\vspace{0.3cm}

\textbf{3. Gustos:}\\[0.1cm]
\es{¿Qué te gusta?}\\
{\scriptsize\color{grau} (Was gefällt dir? $\rightarrow$ Me gusta bailar...)}

\vspace{0.2cm}
\hrule
\vspace{0.2cm}

{\scriptsize\textbf{Hilfe:} \es{Me llamo...} / \es{Soy...} / \es{Me gusta...}}
};
\end{tikzpicture}%
}

% 6 Karten (2 Spalten x 3 Reihen)
\noindent
\begin{tabular}{@{}c@{\hspace{0.5cm}}c@{}}
\karteA & \karteA \\[0.4cm]
\karteA & \karteA \\[0.4cm]
\karteA & \karteA \\
\end{tabular}

\newpage

% ═══════════════════════════════════════════════════════════════════════════
% FRAGEKARTEN B (KOMPLEX) - 6 Stück pro Seite
% ═══════════════════════════════════════════════════════════════════════════

\begin{center}
{\LARGE\bfseries Fragekarten B -- Nivel B1}\\[0.3cm]
{\small 6 Exemplare zum Ausschneiden}
\end{center}

\vspace{0.5cm}

% Karte B - Template
\newcommand{\karteB}{%
\begin{tikzpicture}
\draw[thick, gruppeB, rounded corners=3pt] (0,0) rectangle (8.5,7.5);
\fill[gruppeB] (0,6.8) rectangle (8.5,7.5);
\node[white, font=\bfseries\large] at (4.25,7.15) {TARJETA B -- Nivel B1};

% Inhalt
\node[anchor=north west, text width=7.8cm] at (0.3,6.6) {
\textbf{1. Profesión:}\\[0.05cm]
\es{¿Qué profesión quieres tener? ¿Por qué?}\\
{\scriptsize\color{grau} (Welchen Beruf möchtest du? Warum?)}

\vspace{0.2cm}

\textbf{2. Cualidades:}\\[0.05cm]
\es{¿Qué cualidades son importantes para eso?}\\
{\scriptsize\color{grau} (Welche Eigenschaften sind wichtig dafür?)}

\vspace{0.2cm}

\textbf{3. Tiempo libre:}\\[0.05cm]
\es{¿Qué actividades haces en tu tiempo libre?}\\
{\scriptsize\color{grau} (Was machst du in deiner Freizeit?)}

\vspace{0.2cm}

\textbf{4. Personalidad:}\\[0.05cm]
\es{¿Cómo te describirías en tres palabras?}\\
{\scriptsize\color{grau} (Beschreibe dich in 3 Wörtern.)}

\vspace{0.15cm}
\hrule
\vspace{0.15cm}

{\scriptsize\textbf{Für Zusammenfassung:} \es{Se llama...} / \es{Es...} / \es{Le gusta...} / \es{Quiere ser...}}
};
\end{tikzpicture}%
}

% 6 Karten (2 Spalten x 3 Reihen)
\noindent
\begin{tabular}{@{}c@{\hspace{0.5cm}}c@{}}
\karteB & \karteB \\[0.3cm]
\karteB & \karteB \\[0.3cm]
\karteB & \karteB \\
\end{tabular}

\newpage

% ═══════════════════════════════════════════════════════════════════════════
% ZUSAMMENFASSUNGS-VORLAGE (für Gruppe B) - MIT ANKREUZ-OPTIONEN
% Schneller als freies Schreiben!
% ═══════════════════════════════════════════════════════════════════════════

\begin{center}
{\LARGE\bfseries Zusammenfassungs-Vorlage (Schnell)}\\[0.3cm]
{\small Für Gruppe B -- Ankreuzen + 1 Wort = zeitsparend!}
\end{center}

\vspace{0.5cm}

% Vorlage - Template (kompakter, mit Checkboxen)
\newcommand{\vorlageB}{%
\begin{tikzpicture}
\draw[thick, spanisch, rounded corners=3pt] (0,0) rectangle (8.5,8.5);
\fill[spanisch] (0,7.8) rectangle (8.5,8.5);
\node[white, font=\bfseries\large] at (4.25,8.15) {RESUMEN -- Mi compañero/a};

% Inhalt
\node[anchor=north west, text width=7.8cm] at (0.3,7.6) {
\textbf{Nombre:} \rule{4cm}{0.4pt}

\vspace{0.25cm}

\textbf{1. Carácter} {\scriptsize (máx. 2 ankreuzen)}:\\[0.1cm]
$\square$ \es{simpático/a} \quad $\square$ \es{inteligente} \quad $\square$ \es{creativo/a}\\
$\square$ \es{trabajador/a} \quad $\square$ \es{divertido/a} \quad $\square$ \es{tranquilo/a}

\vspace{0.2cm}

\textbf{2. Gustos} {\scriptsize (máx. 2 ankreuzen)}:\\[0.1cm]
$\square$ \es{bailar} \quad $\square$ \es{leer} \quad $\square$ \es{música}\\
$\square$ \es{deporte} \quad $\square$ \es{viajar} \quad $\square$ \es{cocinar}

\vspace{0.2cm}

\textbf{3. Profesión} {\scriptsize (1 Wort)}:\\[0.1cm]
\es{Quiere ser} \rule{3cm}{0.4pt}

\vspace{0.2cm}

\textbf{4. Por qué} {\scriptsize (1 Wort)}:\\[0.1cm]
\es{porque es} \rule{3cm}{0.4pt}

\vspace{0.15cm}
\hrule
\vspace{0.15cm}

{\scriptsize\textbf{Für Präsentation sagen:}}\\
{\scriptsize \es{Se llama..., es..., le gusta..., quiere ser... porque...}}
};
\end{tikzpicture}%
}

% 4 Karten (2 Spalten x 2 Reihen)
\noindent
\begin{tabular}{@{}c@{\hspace{0.5cm}}c@{}}
\vorlageB & \vorlageB \\[0.3cm]
\vorlageB & \vorlageB \\
\end{tabular}

\newpage

% ═══════════════════════════════════════════════════════════════════════════
% BONUS: TANDEM-ÜBERSICHT (für Lehrkraft)
% ═══════════════════════════════════════════════════════════════════════════

\begin{center}
{\LARGE\bfseries Tandem-Übersicht}\\[0.3cm]
{\small Für die Lehrkraft -- Phase 2 (Gemischte Tandems)}
\end{center}

\vspace{1cm}

\begin{center}
\begin{tabular}{|>{\centering\arraybackslash}p{3cm}|>{\centering\arraybackslash}p{3cm}|>{\centering\arraybackslash}p{5cm}|}
\hline
\cellcolor{gruppeA}\textcolor{white}{\textbf{Gruppe A}} & \cellcolor{gruppeB}\textcolor{white}{\textbf{Gruppe B}} & \textbf{Hinweis} \\
\hline
H.P. (Kl. 10) & E.H. (Kl. 12) & TOP-Tandem: E.H. als Tutor \\
\hline
A.A. (Kl. 10) & V.S. (Kl. 12) & Ruhiges Tandem, klare Struktur \\
\hline
A.S. (Kl. 10) & L.N. (Kl. 11) & A.S. ermutigen! \\
\hline
L.K. (Kl. 11) & C.P. (Kl. 11) & L.K. fokussieren \\
\hline
{[Reserve]} & {[Reserve]} & Falls ungerade Anzahl \\
\hline
\end{tabular}
\end{center}

\vspace{1cm}

\begin{center}
\fbox{%
\begin{minipage}{0.8\textwidth}
\textbf{Ablauf Phase 2:}
\begin{enumerate}
    \item \textbf{B interviewt A} (10 Min) -- B stellt Fragen, A antwortet, B \textbf{kreuzt an}
    \item \textbf{A interviewt B} (5 Min) -- A nutzt einfache Fragen
    \item \textbf{Präsentation} (9 Min) -- \textbf{Nur 3 Tandems!} B stellt A vor (3. Person)
\end{enumerate}

\vspace{0.3cm}

\textbf{Beobachtungsfokus:}\\
\begin{itemize}
    \item Korrekte Verwendung von \es{ser} (Eigenschaft) vs. \es{estar} (Zustand)
    \item Transformation zur 3. Person (\es{es}, \es{le gusta}, \es{quiere})
    \item Kommunikative Kompetenz (Rückfragen, Nachfragen)
\end{itemize}

\vspace{0.3cm}

\textbf{Ausgewählte Tandems für Präsentation:}\\
\begin{enumerate}
    \item E.H. + H.P. (starkes Tandem, zeigt Scaffolding)
    \item L.N. + A.A. (mittleres Tandem)
    \item V.S. + A.S. (zeigt Umgang mit Junge)
\end{enumerate}
\end{minipage}
}
\end{center}

\vspace{1cm}

% Timer-Erinnerung -- ANGEPASSTES TIMING
\begin{center}
\begin{tikzpicture}
\draw[thick, warning, fill=yellow!20] (0,0) rectangle (12,2.5);
\node[font=\Large\bfseries] at (6,1.8) {TIMER nicht vergessen!};
\node[font=\normalsize] at (6,1.1) {Phase 1: \textbf{12 Min} | Phase 2: \textbf{18 Min} | Phase 3: \textbf{9 Min} | Cierre: 3 Min};
\node[font=\small, grau] at (6,0.5) {= 42 Min + 3 Min Puffer = 45 Min};
\end{tikzpicture}
\end{center}

\newpage

% ═══════════════════════════════════════════════════════════════════════════
% VOKABEL-HILFE (Poster für Tafel)
% ═══════════════════════════════════════════════════════════════════════════

\begin{center}
{\LARGE\bfseries Vokabel-Hilfe}\\[0.3cm]
{\small Zum Aufhängen / An die Tafel}
\end{center}

\vspace{0.5cm}

\begin{center}
\begin{tikzpicture}
\draw[very thick, spanisch, rounded corners=5pt] (0,0) rectangle (16,10);
\fill[spanisch] (0,9) rectangle (16,10);
\node[white, font=\bfseries\LARGE] at (8,9.5) {¿CÓMO ES UNA PERSONA?};

% SER-Spalte
\fill[gruppeA!20] (0.3,0.3) rectangle (7.8,8.7);
\node[font=\bfseries\Large, gruppeA] at (4,8.2) {SER + Adjektiv};
\node[font=\large] at (4,7.5) {= Eigenschaft (IMMER so)};

\node[anchor=north, text width=6.5cm, font=\normalsize] at (4,7) {
\begin{tabular}{ll}
\es{simpático/a} & sympathisch \\[0.2cm]
\es{inteligente} & intelligent \\[0.2cm]
\es{trabajador/a} & fleißig \\[0.2cm]
\es{creativo/a} & kreativ \\[0.2cm]
\es{divertido/a} & lustig \\[0.2cm]
\es{tranquilo/a} & ruhig \\[0.2cm]
\es{responsable} & verantwortungsvoll \\[0.2cm]
\es{paciente} & geduldig \\
\end{tabular}
};

\node[font=\itshape, gruppeA] at (4,1.5) {\es{Soy simpático.}};
\node[font=\small] at (4,1) {(Ich bin nett.)};

% ESTAR-Spalte
\fill[gruppeB!20] (8.2,0.3) rectangle (15.7,8.7);
\node[font=\bfseries\Large, gruppeB] at (12,8.2) {ESTAR + Adjektiv};
\node[font=\large] at (12,7.5) {= Zustand (JETZT so)};

\node[anchor=north, text width=6.5cm, font=\normalsize] at (12,7) {
\begin{tabular}{ll}
\es{cansado/a} & müde \\[0.2cm]
\es{enfermo/a} & krank \\[0.2cm]
\es{contento/a} & zufrieden \\[0.2cm]
\es{nervioso/a} & nervös \\[0.2cm]
\es{triste} & traurig \\[0.2cm]
\es{ocupado/a} & beschäftigt \\[0.2cm]
\es{preocupado/a} & besorgt \\[0.2cm]
\es{aburrido/a} & gelangweilt \\
\end{tabular}
};

\node[font=\itshape, gruppeB] at (12,1.5) {\es{Estoy cansado.}};
\node[font=\small] at (12,1) {(Ich bin müde.)};

\end{tikzpicture}
\end{center}

\vspace{1cm}

% Konnektoren
\begin{center}
\begin{tikzpicture}
\draw[thick, spanisch] (0,0) rectangle (16,2.5);
\node[font=\bfseries\large, spanisch] at (8,2) {KONNEKTOREN};
\node[font=\normalsize] at (8,1.2) {
\es{también} (auch) \quad | \quad
\es{además} (außerdem) \quad | \quad
\es{pero} (aber) \quad | \quad
\es{por eso} (deshalb)
};
\node[font=\small, grau] at (8,0.5) {Beispiel: \es{Soy simpático. También soy trabajador. Pero a veces estoy cansado.}};
\end{tikzpicture}
\end{center}

\end{document}
